\section{Preparing a fibration and measuring its relative section}\label{sec:models}

A section regular on the generic fiber may still have poles over the boundary. The following model separates the horizontal and vertical tests for regularity; a smooth resolution will be used only for fiber integrals.

\begin{lemma}[Equidimensional models]\label{val:flat-model}
Let $X\dashrightarrow Y$ be a dominant rational map of integral projective
varieties over $\mathbb C$.  There are a smooth projective birational model
$Y'$ of $Y$ and an integral projective flat family $X_0\to Y'$ with a
birational morphism to the graph of the original map.  Every irreducible
component of every fiber has dimension $\dim X-\dim Y$.  The normalization
$V$ of $X_0$ is equidimensional over $Y'$; in particular, every vertical
prime divisor of $V$ maps onto a prime divisor of $Y'$.  If
$\mathbb C(Y)$ is algebraically closed in $\mathbb C(X)$, the morphism
from $V$, and from any smooth resolution of $V$, to $Y'$ has connected
fibers.  If a torus acts on $X$ and the rational map is invariant, these
constructions and a projective resolution of $V$ can be made equivariant,
with trivial action on the base.
\end{lemma}

\begin{proof}
Embed the projective graph over $Y$ in $\mathbb P^N\times Y$.
Over the flat open set the fibers define a map to the projective
Hilbert scheme \cite[Theorem~3.2 and the Hilbert-scheme specialization
after Proposition~3.8]{Grothendieck1961Hilbert}.
Resolve the closure of its graph, also resolving the base, and pull back the
universal family.  Flatness over an integral base embeds each affine
coordinate ring in its generic localization.  The generic fiber is
integral, so these coordinate rings have neither nilpotents nor zero
divisors.  Thus the pulled-back family is integral and is the closure of
the original family over the flat open set.  Its map to the graph is
birational.

Constancy of the Hilbert polynomial bounds the dimension of every fiber
above by the generic relative dimension.  At any fiber component the
height bound for the ideal generated by local parameters of the smooth
base gives the opposite bound.  This proves purity of fiber dimension.
Normalization is finite, so fiber dimensions and dimensions of images of
components are preserved.  A prime divisor mapping into codimension at
least two in $Y'$ would have dimension at most
$\dim Y'-2+(\dim X-\dim Y')=\dim X-2$, a contradiction.

For connectedness apply Stein factorization to the normal total space.
Its finite intermediate base has function field the algebraic closure
of $\mathbb C(Y')$ in $\mathbb C(X)$; the asserted field condition and
normality of $Y'$ make the finite birational morphism an isomorphism.
The same reasoning applies after resolution.  In the equivariant case,
the family over the original open set is stable under the action.
Its closure and normalization are therefore stable, and equivariant
resolution in characteristic zero supplies the last assertion.
\end{proof}


\subsection{Rank-one adjoints and a continuous maximum metric}\label{har:sec:maximum}

Our first task is to turn a rational relative section into a metrized line on the base. This construction will serve both the character-sensitive transfer and the harmonic applications.
Its normal equidimensional model is used for divisorial normalization
and continuity. A resolution supplies smooth coefficients for the
direct-image theorem; that resolution need not be equidimensional.

\begin{proposition}[Divisorial normalization and maximum metric]\label{har:prop:maximum}
Let \(X,Y\) be smooth connected projective, let \(L=-K_X\) have a smooth
semipositive metric, and suppose
\[
 V\xrightarrow{\nu}W\xrightarrow{\mu}X,\qquad g:W\to Y,
 \qquad f=g\nu,
\]
where \(W\) is normal, \(g\) is projective surjective equidimensional,
\(\mu\) is birational, and \(V\) is a smooth projective resolution.
Set \(D_V=(\mu\nu)^*L\) and \(E=K_{V/X}\), and let \(s_E\) be the canonical Jacobian section of \(E\).
Suppose a nonzero rational section \(s\) of \(b\mu^*L\), \(b>0\), has
no horizontal poles and, on the generic fiber $V_\eta$, satisfies
\[
 H^0(V_\eta,E+kbD_V)=\C(Y)\,s_Es^k
 \qquad(k\ge0).
\]
Alternatively, suppose a torus acts compatibly on the diagram, trivially
on $Y$, and let $P$ be $L$ with any chosen linearization. Put
$P_V=(\mu\nu)^*P$. In this case it suffices that $s$ be an invariant
rational section of $b\mu^*P$ and that the displayed condition be
replaced by
\[
 H^0(V_\eta,E+kbP_V)^T=\C(Y)\,s_Es^k
 \qquad(k\ge0),
\]
where $E$ carries its natural Jacobian action.

After replacing \(b,s\) by a multiple and a power, there are a
line \(Q\) on \(Y\) and a regular injection
\[
 g^*Q\longrightarrow b\mu^*L
\]
whose squared fiber-maximum norms define a positive continuous metric
with plurisubharmonic weights on \(Q\).
In the invariant case $Q$ has trivial linearization and the injection
$g^*Q\to b\mu^*P$ is equivariant.
\end{proposition}

\begin{proof}
In the invariant case average a smooth semipositive metric on $L$ over
the maximal compact torus. The chosen linearization of $P$ differs
from the natural one by a character, which is unitary on this compact
torus, so the averaged metric is invariant for both actions.

The rank assumption at $k=0$ forces the generic fiber to be geometrically connected. Indeed its constant functions, multiplied by $s_E$, inject into that adjoint space. In the invariant case the connected torus acts trivially on the finite Stein base, so all these constants are invariant. Their dimension is therefore one. Smooth general fibers are consequently connected and irreducible.

For every prime divisor \(B\) of \(Y\), put
\[
 c_B=\min_{\Gamma\mapsto B}
       \frac{\operatorname{ord}_{\Gamma}(s)}
            {\operatorname{ord}_{\Gamma}(g^*B)}.
\]
Properness and equidimensionality supply divisorial lifts of \(B\);
all vertical prime divisors lie over base primes. Only finitely many
\(c_B\) are nonzero. Clear denominators and use the resulting integral
divisor as \(Q\). Sending its rational canonical section to the powered
\(s\) gives a regular bundle map by normality. Its effective Cartier
zero divisor \(G\) omits at least one component above every base prime.
In the invariant case all these divisors are invariant, the base line carries trivial linearization, and the map is to \(b\mu^*P\). In an eigensection version, one
instead twists the base linearization by the section character.

For a local frame \(e\) of \(Q\), denote its image by \(s_e\) and put
\begin{equation}\label{har:eq:max}
 M_e(y)=\max_{x\in W_y}|s_e(x)|^2.
\end{equation}
This is continuous. Properness gives upper semicontinuity. For lower
semicontinuity use that an equidimensional holomorphic map from a pure
dimensional normal space to a smooth base is open. To verify openness
at a point, cut the source by as many general local hyperplanes as the
fiber dimension, isolating the point in its fiber. The slice germ has
dimension at least that of the base and is finite over that smooth
base germ. Its image is therefore the entire base germ. This can be
done inside any prescribed source neighborhood.

The locus \(Z_0\) where \(M_e\) vanishes is algebraic closed: its
complement is the constructible image \(g(W\setminus\operatorname{Supp}G)\),
and continuity makes \(Z_0\) closed in the classical topology; a constructible
classically closed subset is Zariski closed. By the omitted-lift
property, \(Z_0\) has codimension at least two.

We show first that \(-\log M_e\) is plurisubharmonic on a dense open.
Work on a fixed smooth-fibration open for \(f\), where the canonical
exceptional section and \(s_e\) are not identically zero on any fiber
component. Take a local frame \(\eta\) of \(K_Y\). The section
\[
  (s_E/f^*\eta)s_e^k
 \quad\text{of }K_{V/Y}+(kb+1)D_V
\]
generates the rank-one adjoint space at the generic point. In the shifted equivariant case its coefficient is \(D_V+kbP_V\); as a metrized ordinary line it is still \((kb+1)D_V\). Berndtsson's
smooth direct-image theorem \cite[Theorem~1.2]{Berndtsson2009} applies on
the bundle and base-change locus with smooth semipositive coefficient
\((kb+1)D_V\). In the invariant case the invariant line is an
\emph{orthogonal quotient} of the full direct image: the compact torus
acts unitarily, and the sum of its nonzero-weight spaces is its
orthogonal complement. It therefore has semipositive curvature.
This quotient argument, rather than positivity of arbitrary sublines,
is what permits taking invariants.

For $k\ge1$ the resulting plurisubharmonic function is
\begin{equation}\label{har:eq:moments}
 -\frac1k\log\int_{V_y}|s_e|^{2k}\,d\mu_y,
\end{equation}
where \(\mu_y\) comes from \(s_E/f^*\eta\) with the metric of \(D_V\).
It is smooth, nonnegative, of full support, and independent of \(k\).
Each integral is smooth positive on our fixed open, so the curvature
inequality extends there from its possibly \(k\)-dependent base-change
open. Normalized to probability measures, its \(1/k\)-powers increase
to \(M_e\). They converge locally uniformly: a point just below the
positive fiber maximum has a neighborhood of uniformly positive
\(\mu_y\)-mass on nearby fibers; compactness supplies finitely many such
neighborhoods, while bounded total mass gives the upper estimate.
Thus \eqref{har:eq:moments} converges locally uniformly to \(-\log M_e\).

By continuity this weight extends plurisubharmonically across the bad
analytic loci wherever \(M_e>0\). It cannot tend to \(+\infty\) across
\(Z_0\) of codimension at least two. Through a point of $Z_0$ choose
a small analytic disc whose boundary misses $Z_0$. Nearby parallel
discs have their boundaries in a fixed compact set off $Z_0$.
Since $Z_0$ has codimension at least two, a dense set of these discs
misses it entirely. The maximum principle therefore bounds the weight
uniformly at their centers, and continuity off $Z_0$ supplies the
same bound there. A zero of the continuous \(M_e\) would contradict
this bound. Hence \(M_e>0\) everywhere, and the weight is continuous
plurisubharmonic. Its transformation rule under a change of frame is
exactly the metric transformation rule for \(Q\).
\end{proof}
